\begin{center}
{\heiti\zihao{3} 帕金森病的特征识别和声音诊断模型}
\end{center}

\vspace{1em}

\noindent{\heiti 摘要：}
帕金森病常伴随发声稳定性、声带振动和谱结构等方面的改变，语音声学特征为辅助筛查和探索性分型提供了低侵入、可重复采集的数据基础。本文围绕两个含重复录音结构的帕金森病语音数据集，建立受试者级二分类辅助诊断模型、判别性声学特征综合评价模型以及 PD 内部潜在语音表型无监督分型模型。针对同一受试者多次录音可能导致的样本单位错误，本文将受试者作为核心分析单位；在监督诊断模型中采用按受试者分组的交叉验证，并将标准化、模型训练和相关特征处理限制在训练折内部完成。

对问题一和问题三，本文构建并比较 Logistic Regression、SVM-RBF 和 Random Forest 三类二分类辅助诊断模型。Dataset 1 主分析采用受试者级指标，Random Forest 的准确率和 AUC 分别为 0.8290 和 0.8511，但其特异度为 0.4385，显示模型更偏向识别 PD 类；SVM-RBF 的 balanced accuracy 为 0.7118，综合表现较稳健。Dataset 2 按受试者 ID 分组交叉验证，SVM-RBF 达到 accuracy 0.8375、balanced accuracy 0.8375、AUC 0.8906、specificity 0.8750，作为该数据集的主推荐模型。

对问题二，本文在 Dataset 1 受试者级数据上整合 Mann-Whitney U 检验、Welch t-test、FDR-BH 校正、Cohen's \(d\)、Cliff's delta、L1 Logistic 稳定选择、ExtraTrees 重要性和 held-out permutation importance，得到候选语音生物标志物综合排序。Top20 特征主要包括 Delta / Delta-delta 动态谱波动特征与 TQWT 多尺度特征；Top50 中 TQWT 特征 31 个、Delta / Delta-delta 特征 16 个、Shimmer、MFCC 和 Nonlinear 特征各 1 个。上述结果仅表示判别性声学特征或候选语音生物标志物，不构成临床因果解释。

对问题四和问题五，由于数据集中没有真实运动型/非运动型标签，也没有静止性震颤、运动迟缓、肌肉僵直、疼痛、痴呆、睡眠障碍六类临床症状标签，本文仅在 Dataset 1 的 188 名 PD 受试者内部进行无监督潜在语音表型分析。二类分型主方案为 \feature{top20_biomarkers} / \feature{pca_5} / KMeans，得到两簇人数 62 和 126，silhouette 为 0.3674，stability ARI mean 为 0.9854。六类分型主方案为 \feature{top20_biomarkers} / \feature{pca_5} / AgglomerativeClustering，六簇人数为 13、55、34、28、25、33，silhouette 为 0.1823，stability ARI mean 为 0.4582。六类方案的聚类质量弱于二类方案，因此只作为探索性六类潜在语音表型，需要真实临床症状标签进一步验证。

本文在建模流程中保持受试者分组、训练折内标准化、主声学模型排除非声学变量和无监督标签不回流输入等约束，以降低样本单位错误和信息泄漏风险。综上，语音声学特征对帕金森病辅助筛查、判别性特征识别和探索性潜在语音表型分析具有一定价值，但本文模型不应解释为临床确诊工具，仍需更大样本和真实临床症状标签验证。

\vspace{1em}
\noindent{\heiti 关键词：} 帕金森病；语音特征；受试者级交叉验证；候选语音生物标志物；潜在语音表型；无监督聚类
